% 第 21 章。数字：checks/check_ch21.py；引文：data/snapshots/divergence_facts.csv（逐字核对见
% data/fetch/verify_web_quotes.py）；数据：maddison2023_divergence、gpih_allen2011_welfare、
% clark2015_england_decadal_1250_1860、boe_millennium_england_1270_1913；图：figsrc/fig_ch21.py。
\chapter{工业革命与大分流}\label{ch:divergence}

\begin{chapterguide}
\guidequestion{人均收入几千年几乎不增长，为什么突然起飞？工业化为什么先在英国，而不是在江南发生？}
\begin{guidegoals}
\item 用马尔萨斯模型（1798）解释前工业时代人均收入的停滞；
\item 用 Maddison 数据描述起飞的时间与幅度，并说明数据的不确定性；
\item 比较 Pomeranz (2000) 的煤炭与殖民地解释、Allen (2009) 的高工资解释与 Mokyr 的知识解释；
\item 评估各种解释的证据，以及它们对今天发展问题的启示。
\end{guidegoals}
\guideprereq{\chref{ch:medieval}\enspace\nameref{ch:medieval}；\chref{ch:accounts}\enspace\nameref{ch:accounts}}
\end{chapterguide}

1798 年，马尔萨斯（Thomas Robert Malthus，1766—1834）匿名出版了《人口原理》（\emph{An Essay on the Principle of Population}）的第一版\cite{malthus1798}。论证从两条公设出发：“第一，食物是人类生存所必需的”；“第二，两性间的情欲是必然的，并将几乎保持现状”\cite{malthus1798}。由此，“人口增长的能力无限大于土地为人类生产生活资料的能力”：“人口在不受抑制时按几何级数增长，生活资料只按算术级数增长。”\cite{malthus1798}北美的人口“二十五年就翻了一番”，他就把这定为规则：人口在不受抑制时“每二十五年翻一番”\cite{malthus1798}。

他拿英国来算。“这个岛的人口据计算约为七百万”；人口每二十五年翻一番，生活资料每二十五年增加够七百万人吃的一份，“第一个世纪结束时，人口将是一亿一千二百万，而生活资料只够养活三千五百万，剩下七千七百万人完全没有着落”\cite{malthus1798}。这当然不会真的发生：人口远在超过食物之前，就会被贫困与晚婚压回来。所以他对穷人的前途是悲观的：“国家财富的增长对改善劳动贫民的处境几乎没有或完全没有作用。”\cite{malthus1798}

马尔萨斯写书的时候，英格兰正在走出他描述的那个世界。按英格兰银行“千年宏观数据”里 Broadberry 等人的估计，英格兰的人口 1800 年是 862 万，1850 年是 1,651 万，增加到 1.915 倍\cite{boe2017,broadberry2015}。这没有马尔萨斯说的不受抑制时那么快——按每二十五年翻一番，五十年后应是 3,448 万——却比英格兰以往任何一个时期都快。人均实际 GDP 也没有被增加的人口吃掉，反而增加了 25.0\%。把时间拉长：从 1270 年到 1700 年，英格兰的人均实际 GDP 增加到 2.07 倍，用了四百多年；从 1700 年到 1913 年，增加到 3.26 倍\cite{boe2017,broadberry2015}。别的地方没有跟上。按 Maddison 项目数据库，英国的人均 GDP 1700 年是中国的 1.56 倍，1913 年是 8.34 倍\cite{mpd2023}。这个差距的扩大，Pomeranz（2000）用作书名，叫作\textbf{大分流}（great divergence）\cite{pomeranz2000}。

这一章回答导引框里的两个问题：人均收入几千年几乎不增长，为什么突然起飞？工业化为什么先在英国，而不是在江南发生？依次讨论：马尔萨斯写书时的英格兰、中国与世界（\secref{sec:divergence-world}）；马尔萨斯的模型，以及走出它的三条路（\secref{sec:divergence-malthus}）；起飞的时间、幅度与数字的可靠性（\secref{sec:divergence-takeoff}）；分流从什么时候开始（\secref{sec:divergence-when}）；为什么是英国（\secref{sec:divergence-why}）。\secref{sec:divergence-test} 用英格兰 1250—1860 年的数据检验马尔萨斯模型。本章引用的外文史料，中文都是本书的翻译。Pomeranz（2000）、Allen（2009）、Mokyr（2009）与 Clark（2007）四本书没能取得，书中的论点据 EH.Net 的书评转述；正文注明哪些是书评人的概括，哪些是书评引的原文。

\section{当时的世界}\label{sec:divergence-world}

这一节交代：马尔萨斯写书的时候，英格兰的人口怎样变化；他与斯密怎样看中国；世界各地的人均收入相差多少。

\subsubsection{英格兰：人口在加速}
马尔萨斯说，欧洲的主要国家现在人口增长很慢，“不是每二十五年翻一番，而是要三四百年或更久”\cite{malthus1798}。\chref{ch:smith}的末尾算过，按千年数据，英格兰的人口 1700—1770 年平均约 205 年翻一番，1750—1780 年已经加快到约 107 年。加速没有停：1780—1820 年平均 61.0 年翻一番，1800—1850 年 53.3 年\cite{boe2017,broadberry2015}。马尔萨斯写书，正赶上英格兰人口有史以来增长最快的几十年。

他也看到了财富的增长。他承认，过去一个世纪英国的商业“确实一直在迅速前进”，但增长“主要是劳动的产品，而不是土地的产品”，所以“维持劳动的实际基金增加得很慢”\cite{malthus1798}。至于劳动贫民，“我相信，他们并没有更多地支配生活必需品与便利品”；“比起革命时期，他们中有大得多的一部分在制造业里做工，挤在狭窄而不卫生的屋子里”\cite{malthus1798}。革命时期指 1688 年的光荣革命。马尔萨斯的判断，财富增长而劳动者的处境不变，\secref{sec:divergence-test} 会用工资的数据检验。

\subsubsection{斯密与马尔萨斯眼中的中国}
斯密说中国停滞（\secref{sec:smith-states}）。马尔萨斯在第十六章转述他：“亚当·斯密博士说，中国大概早已富到它的法律与制度所容许的程度”，又说“换一套法律与制度，再重视对外贸易，它还可以富得多”\cite{malthus1798}。马尔萨斯接受当时的说法，“中国是世界上土地最肥沃的国家，几乎所有的土地都在耕种”\cite{malthus1798}。他关心的是，财富的增加能不能改善中国下层人民的处境。他的回答是不能。中国“相对于它的资本所能雇用的人数，人口相当过多”，“劳动如此充裕，没有人费心去节省它”\cite{malthus1798}。如果对外贸易受到重视，中国“靠劳动者之多、劳动之便宜，可以制造数量极大的制成品向外国出售”，换回来的却“主要是从世界各地收集来的奢侈品”，不是能养活更多人的粮食\cite{malthus1798}。这两句话后来分别成了两场讨论的起点：劳动便宜的地方，没有人去发明节省劳动的机器（\secref{sec:divergence-why}）；靠廉价劳动制造出口品，是 20 世纪东亚的道路（\chref{ch:eastasia}）。

\subsubsection{世界：六个经济体}
图~\ref{fig:div_world}(a) 画出了 Maddison 项目数据库 2023 年版里六个经济体 1500—1913 年的人均 GDP（按 2011 年国际元计）\cite{mpd2023,boltvanzanden2024}。数据库里各国的序列来自不同的研究，图注列出了它们\cite{broadberry2015,vanzanden2012,smits2000,bgl2018,xu2016,bcg2015,bassino2018,fukao2015,mccusker2006,sutch2006,prados2009}。1700 年，荷兰最富，3,377 元；英国 2,412 元；中国 1,543 元，印度 1,162 元，日本 1,073 元。1820 年，英国 3,306 元，荷兰降到 3,006 元，美国 2,674 元，日本 1,317 元，中国降到 882 元。1913 年，美国 10,108 元，英国 8,212 元，荷兰 6,454 元，日本 2,431 元，印度 1,073 元，中国 985 元。英国在 1808 年超过荷兰（数据库里荷兰 1809—1814 年没有数字），此后一直领先；美国在 1880 年第一次超过英国，此后两国互有高下，1899 年以后美国一直领先。图~\ref{fig:div_world}(b) 是英格兰自己的长时段：黑死病以后人口减少、人均产出上升，两者此消彼长；1700 年以后，两者一起上升。这些数字的来历与可靠性，是\secref{sec:divergence-takeoff}与\secref{sec:divergence-when}的问题。

\fig{div_world}{人均 GDP 的分流，1270—1913 年。(a) 看各条线什么时候分开：六个经济体的人均 GDP，按 2011 年国际元计，纵轴为对数刻度；1820 年以前多数经济体只有若干基准年份的估计，点之间用直线连起。英国一栏 1700 年以前只是英格兰。(b) 看两条线什么时候不再此消彼长：英格兰的人口（1270—1870）与人均实际 GDP（1270—1913，按 2013 年价格），都以 1270 年 $=100$，纵轴为对数刻度；浅灰阴影是黑死病（1348—1350）。数据：(a) Maddison 项目数据库 2023 年版（Bolt 与 van Zanden 2024），各国的原始来源为：英国，Broadberry 等（2015）；荷兰，van Zanden 与 van Leeuwen（2012）、Smits、Horlings 与 van Zanden（2000）；中国，Broadberry、Guan 与 Li（2018）、Xu 等（2016）；印度，Broadberry、Custodis 与 Gupta（2015）；日本，Bassino 等（2018）、Fukao 等（2015）；美国，McCusker（2006）、Sutch（2006）、Prados de la Escosura（2009）。(b) 英格兰银行《千年宏观数据》表 A2 与 A21（Broadberry 等 2015）。}{fig:div_world}

\section{马尔萨斯的模型与三条出路}\label{sec:divergence-malthus}

这一节回答：马尔萨斯的模型说了什么？在什么条件下，人均收入可以持续增长？

\subsubsection{两种抑制}
马尔萨斯描述的是一个往复的过程。人口增加了，“原来养活七百万人的食物，现在要分给七百五十万或八百万人”；劳动者多于市场上的工作，“劳动的价格必定趋于下降，同时食物的价格趋于上涨”，“劳动者于是必须更辛苦地劳动，才能挣到和以前一样多的东西”\cite{malthus1798}。在这个困苦的时期，“结婚的阻碍与养家的困难如此之大，人口停滞不前”；同时，劳动便宜了，又“鼓励耕种者在土地上使用更多的劳动，开垦新的土地，给已经耕种的土地更充分地施肥与改良”，生活资料慢慢追上人口；“劳动者的处境于是又相当舒适了，对人口的限制在某种程度上放松”，“幸福方面同样的倒退与前进的运动周而复始”\cite{malthus1798}。人口靠两种抑制压回来：“对养家之难的预见起着预防性抑制的作用；下层阶级中一些人的实际困苦，使他们不能给孩子适当的食物与照料，则起着积极性抑制的作用。”\cite{malthus1798}

\subsubsection{写成式子}
\chref{ch:medieval}用 Ashraf 与 Galor 的模型写过马尔萨斯的经济\cite{ashrafgalor2011}。产出 $Y=(AX)^{\alpha}L^{1-\alpha}$（式~\eqref{eq:medieval-output}），$X$ 是土地，$L$ 是人口，$A$ 是技术，$\alpha$ 是土地在产出里的份额；每人的收入 $y=(AX/L)^{\alpha}$；收入高于稳态 $\bar y$，人口就增加，$L_{t+1}/L_t=y_t/\bar y$。结论是式~\eqref{eq:medieval-dynamics}：一期是一代人，收入对稳态的偏离每过一代只剩 $1-\alpha$ 次方；技术与土地长期只决定人口，不决定收入。

把这个模型推广两步。第一，技术每代都在进步。土地不变时，$\ln y_{t+1}-\ln y_t=\alpha(\ln A_{t+1}-\ln A_t)-\alpha(\ln L_{t+1}-\ln L_t)$；记 $g=\alpha(\ln A_{t+1}-\ln A_t)$，它是人口不变时技术每代使收入增长的幅度，也就是全要素生产率 $(AX)^{\alpha}$ 每代的增长率。第二，人口对收入的反应可强可弱：$L_{t+1}/L_t=(y_t/\bar y)^{\beta}$，$\beta=1$ 是\chref{ch:medieval}的情形，$\beta$ 越小，收入高出稳态时人口增加得越少。记 $z_t=\ln(y_t/\bar y)$ 为收入对稳态的偏离（取对数），两步合起来：
\begin{keyeq}
\ln\frac{y_{t+1}}{y_t}=g-\alpha\ln\frac{L_{t+1}}{L_t},\quad \frac{L_{t+1}}{L_t}=\Big(\frac{y_t}{\bar y}\Big)^{\beta}\qquad\Longrightarrow\qquad z_{t+1}=(1-\alpha\beta)\,z_t+g,\qquad z^{*}=\frac{g}{\alpha\beta} .\label{eq:divergence-escape}
\end{keyeq}
$g=0$、$\beta=1$ 时，$z_{t+1}=(1-\alpha)z_t$，就是式~\eqref{eq:medieval-dynamics} 取对数。

式~\eqref{eq:divergence-escape} 说，只要 $\alpha\beta>0$，收入就趋向一个平台 $z^{*}$。技术进步得越快（$g$ 越大），平台越高；可是到了平台，收入就不再增长，技术的每一步进步都变成人口：人口每代增长 $\beta z^{*}=g/\alpha$，收入停在原处。马尔萨斯的经济并不是没有技术进步，而是技术进步表现为人更多，而不是人更富。

走出这个陷阱有三条路，各对应后面的一类解释。第一，$g$ 提高：技术进步加快，平台升高，向平台移动的那些代里收入在增长；Mokyr 的知识、Allen 的诱致性发明，回答的是 $g$ 为什么在英国提高。第二，$\alpha$ 下降：土地在生产里不那么重要了，煤代替了木柴，新大陆的土地补上了本国的土地，工业与服务业的份额超过了农业；Pomeranz 的煤与殖民地属于这一条。第三，$\beta$ 下降：收入提高以后，人口不再同比例地增加，人们晚婚、少生，即马尔萨斯的预防性抑制占了上风；这是 Clark 的论题的一部分，也是后来的人口转变（\chref{ch:demography}）。$\alpha\beta\to0$ 时，$z_{t+1}=z_t+g$，收入每代增长 $g$，这就是现代经济增长。

还有一个更直接的界限。人口增长有生理的上限：马尔萨斯说不受抑制时每二十五年翻一番，即每年约 2.77\%。由式~\eqref{eq:divergence-escape} 的第一式，按年计，收入每年至少增长 $g_{A}-\alpha n_{\max}$，$g_{A}$ 是每年的全要素生产率增长，$n_{\max}$ 是人口增长的上限。只要全要素生产率增长快过 $\alpha n_{\max}$，人口增加得再快，也吃不掉全部的进步。土地份额 $\alpha=0.25$ 时，这个门槛是每年 0.69\%；$\alpha=0.15$ 时，是 0.42\%。

\begin{example}{平台有多高}\label{ex:divergence-plateau}
一代按 30 年计。（1）设 $\alpha=0.25$，$\beta=1$，全要素生产率每年增长 0.25\%。按式~\eqref{eq:divergence-escape}，平台 $z^{*}$ 是多少？收入比稳态 $\bar y$ 高多少？到了平台，人口每年增长多少？（2）从 $z_0=0$ 出发，一代以后收入高出稳态多少？离平台的差距（取对数）多少代减半？（3）设 $\alpha=0.15$，全要素生产率每年增长 0.8\%，$\beta$ 仍是 1。平台、差距减半的代数与到达平台时的人口增长各是多少？这样的人口增长可能吗？
\solution （1）$g=30\times0.25\%=0.075$，$z^{*}=0.075/(0.25\times1)=0.3$，收入是稳态的 $e^{0.3}=1.35$ 倍，高 35\%。人口每代增长 $g/\alpha=0.3$，每年 1\%。技术每年进步 0.25\%，收入却只比稳态高 35\%，其余的都变成了人口。（2）$z_1=0.075$，收入高出 7.8\%。差距每代乘以 $1-\alpha\beta=0.75$，减半需要 $\ln0.5/\ln0.75=2.41$ 代，约 72 年。（3）$g=0.24$，$z^{*}=0.24/0.15=1.6$，平台是稳态的 $e^{1.6}=4.95$ 倍。差距每代乘以 0.85，减半要 $\ln0.5/\ln0.85=4.27$ 代，约 128 年。一代以后收入高出 $e^{0.24}-1=27\%$；四代（120 年）以后，$z_4=1.6\times(1-0.85^{4})=0.765$，收入是稳态的 2.15 倍。到达平台时，人口要每代增长 $g/\alpha=1.6$，每年 5.3\%，远远超过马尔萨斯说的不受抑制的速度（每年 2.77\%）。所以，即使人口对收入的反应不变，土地份额低、技术进步快的经济也会有一个多世纪的收入增长；人口根本跟不上，收入每年至少增长 $0.8\%-0.15\times2.77\%=0.38\%$。
\end{example}

\begin{intuition}[跑步机]
马尔萨斯的经济像一台跑步机。技术进步让人往前跑，人口增长让传送带往后转；人跑得越快，传送带也转得越快，最后人停在原地，只是跑步机上站的人多了。式~\eqref{eq:divergence-escape} 说，要离开原地，或者跑得比传送带能达到的最快速度还快（$g_{A}>\alpha n_{\max}$），或者让传送带对人的速度不那么敏感（$\beta$ 变小），或者让传送带的作用变弱：土地在生产里的份额 $\alpha$ 变小，人多了，每个人分到的地少了，也不那么要紧。
\end{intuition}

\fig{div_model}{走出马尔萨斯陷阱的三条路（示意图）。(a) 看平台的高低与到达平台要多久：按式~\eqref{eq:divergence-escape}，从稳态（$z_0=0$）出发十代（一代 30 年）的人均收入，以稳态 $\bar y=1$；实线是例~\ref{ex:divergence-plateau} (1) 的情形（$\alpha=0.25$，$\beta=1$，全要素生产率每年增长 0.25\%），褐色虚线把 $\beta$ 减半，浅色线是例~\ref{ex:divergence-plateau} (3) 的情形（$\alpha=0.15$，每年 0.8\%）。(b) 看平台从哪里来：下一代的 $z_{t+1}$ 对这一代的 $z_t$，两条线与 45 度线的交点是平台 $z^{*}$，阶梯是实线情形的前几代。参数是教学假设。}{fig:div_model}

\begin{compare}{马尔萨斯经济与现代经济增长}{马尔萨斯经济}{现代经济增长}
\cmprow{技术进步变成什么}{更多的人口；人均收入停在平台上}{更高的人均收入}
\cmprow{土地的份额 $\alpha$}{大：英格兰 1250—1630 年平均 0.235}{小：1860 年代 0.095}
\cmprow{人口对收入的反应 $\beta$}{强：收入高了就早婚、多生、少死}{弱：人口增长与收入脱钩，后来收入越高生育越少（\chref{ch:demography}）}
\cmprow{全要素生产率}{慢：英格兰 1250—1640 年每年 $-0.04\%$}{快：1800—1860 年每年 0.80\%}
\cmprow{人均收入与人口}{反向变动：人多了，工资低}{同向变动：两者一起增长}
\cmprow{本章的检验}{1250—1630 年，人均收入对人口的弹性 $-0.72$}{1800—1860 年，弹性 $+0.46$}
\end{compare}

\section{起飞：时间、幅度与数字的可靠性}\label{sec:divergence-takeoff}

这一节回答：英国的人均收入什么时候开始持续增长？增长有多快？这些数字有多可靠？

\subsubsection{慢起飞}
在 Maddison 数据库里，英国的人均 GDP 1700—1820 年每年只增长 0.26\%，1820—1870 年加快到 1.13\%，1870—1913 年是 0.80\%\cite{mpd2023,broadberry2015}。按每年 0.26\% 的速度，翻一番要约 264 年。同一时期，荷兰每年下降 0.10\%，中国每年下降 0.47\%，日本每年增长 0.17\%。1820—1870 年，美国每年增长 1.17\%，1870—1913 年 1.73\%；日本 1870—1913 年 1.00\%\cite{mpd2023}。工业革命在英国不是一夜之间的跳跃：18 世纪的人均增长很慢，快起来是在 1820 年以后；可它比以往任何时代的增长都持续。

增长集中在少数部门，最突出的是棉纺织业。Broadberry 与 Gupta 的表 1 给出了英国原棉的年消费量：1697—1699 年与 1700—1709 年都是 110 万磅，1770—1779 年 480 万磅，1780—1789 年 1,550 万磅，1790—1799 年 2,860 万磅，1800—1809 年 5,960 万磅，1820—1829 年 1.665 亿磅，1840—1849 年 5.263 亿磅\cite{bg2005}。“1780 年到 1800 年，产出每年增长 10.8\%，出口则以惊人的每年 14.0\% 扩张。”\cite{bg2005}机器是原因。按 Catling 的估计，1780 年前后，在印度把 100 磅棉花纺成 80 支纱大约要 50,000 个工时；同年在英国，用克朗普顿（Samuel Crompton，1753—1827）的走锭纺纱机要 2,000 工时，1825 年用罗伯茨（Richard Roberts，1789—1864）的自动走锭纺纱机只要 135 工时\cite{bg2005}。

\subsubsection{数字有多可靠}
1820 年以前的人均 GDP 都是估计出来的。Broadberry、Guan 与 Li 给他们的中国序列打了等级：“对于可靠的数字（A 级），所报告的序列的误差范围被判断为不到正负 5\%”；“对于较好的估计（B 级），误差范围为正负 5\%—15\%”；“粗略的估计（C 级）为正负 15\%—25\%，猜测（D 级）则超过正负 25\%”\cite{bgl2017}。按他们的表 3，人均 GDP 在北宋与明代是 B 级，在清代是 A 级；名义 GDP 在北宋与明代是 C 级，在清代是 B 级\cite{bgl2017}。人口也有误差。Perkins（1969）给他的基准年份估计了一个范围，按 Broadberry 等人的转述，这个范围“从 1393 年的约正负 10\% 降到 19 世纪中叶的正负 6\%”；他们说，“也许更恰当的是采纳 Perkins 的建议，即 80\% 的置信区间”\cite{bgl2017}。

英格兰的数字也有两套。Broadberry 等人从产出一方估计，把农业、工业与服务业的产出加起来；Clark 从收入一方估计，把工资、地租与资本的收入加起来\cite{broadberry2015,clark2010}。两套数字 1700 年以后大体一致，在此以前相差很大。Broadberry 等人的十年平均以 1860 年代 $=100$，Clark 的指数照他的工作簿（1860 年代是 100.3）：1270 年代，Broadberry 等人的人均实际 GDP 是 23.0，Clark 的人均实际国民净收入是 44.8；黑死病以后的 1450 年代，前者 31.9，后者 91.1；1640 年代，30.8 与 52.5；1700 年代，50.0 与 59.0；1800 年代，66.3 与 69.6；1850 年代，90.4 与 94.9（图~\ref{fig:div_england}(b)）。按 Broadberry 等人，1270—1700 年人均收入增加到 2.07 倍，中世纪的英格兰已经在缓慢地增长；按 Clark，1450 年代的人均收入几乎赶上 1850 年代，此后又跌回去，前工业时代的英格兰是一个典型的马尔萨斯经济。哪一套对，决定了我们怎样理解起飞：是在一条缓慢上升的路上加速，还是从一条水平线上起飞。

\section{分流从什么时候开始}\label{sec:divergence-when}

这一节回答：18 世纪的长江三角洲与英格兰，谁的生活水平更高？分流是在工业革命之前还是之后开始的？

\subsubsection{Pomeranz：1750 年的相似}
Lal 在书评里概括 Pomeranz 的结论：他“相当令人信服地表明，直到十八世纪的世纪之交，中国与欧洲的核心地区之间生活水平并没有明显的分流”；“晚至 1750 年，长江三角洲与英格兰之间的相似之处仍多于不同之处”\cite{pomeranz2000,lal2000}。“核心地区”是关键：Pomeranz 比较的不是中国与欧洲，而是两者最发达的地区，长江三角洲（江南）与英格兰。问题于是变成，按 Lal 的转述：“那么，英格兰与随后的欧洲为什么没有走它们远东表亲那种劳动密集的‘勤劳革命’之路，而走上了资本密集的工业革命之路？”\cite{lal2000}

\subsubsection{工资：福利比率}
Allen、Bassino、Ma、Moll-Murata 与 van Zanden 用工资来比较\cite{allen2009wp,allen2011china}。他们先定一个成年人维持“仅够糊口”（bare bones）生活一年所需的一篮子物品，“乘以三，代表一个家庭的年度预算”，“我们现在再把它们提高 5\%，这是对房租最低限度的估计”；“我们假定一年的工作是 250 天”\cite{allen2009wp}。一个非熟练劳动者全年的收入除以这个家庭预算，就是他们的指标：“这类实际工资指数叫作‘福利比率’。”比率“等于一，一个全时工作的非熟练劳动者挣的正好够按糊口的收入养活一家”\cite{allen2009wp}。

图~\ref{fig:div_wages}(a) 画出了他们的福利比率\cite{allen2011china}。1750 年代，伦敦是 3.96，阿姆斯特丹 4.18；米兰 1.34，北京 1.35，广州 1.23，苏州 1.16；京都 0.81，孟加拉 0.88。他们的结论是，在他们比较的欧洲与中国的城市里，“中国的城市与意大利的城市并列最后一名，后者是欧洲生活水平最低的”；长江三角洲号称中国最发达的地区，“那里的实际工资并不明显高于北京或广州”\cite{allen2009wp}。18 世纪的大部分时间里，“中国与日本主要城市的非熟练劳动者——尽管很穷——同中欧与南欧的同类劳动者生活水平大体相同”，“只有英格兰与低地国家领先于其他地区”\cite{allen2009wp}。所以他们“反驳修正派的说法，即长江三角洲这样的中国先进地区在工业革命前夕与英格兰不相上下”\cite{allen2009wp}。

此后，北京的实际工资“继续下滑，直到世纪中叶的太平天国起义，那时达到了危及生命的低水平”；“秩序恢复以后，生活水平缓慢改善，一直到 20 世纪初”\cite{allen2009wp}。按他们的数据，北京的福利比率 1840 年代平均 0.76，1850 年代 0.56，1900—1913 年 1.49，差不多回到 1750 年代的水平；伦敦则从 1750 年代的 3.96 升到 1850 年代的 4.53、1900—1913 年的 7.77。所以他们说：“1913 年欧洲与中国之间的差距，大部分来自欧洲的进步，而不是中国的衰退。”\cite{allen2009wp}阿姆斯特丹是一个反例：它的福利比率从 1750 年代的 4.18 降到 1850 年代的 2.31；Allen 等人说，“荷兰经济在 19 世纪初步履蹒跚”，“不过，到世纪中叶，增长恢复了，实际工资攀上新的高度”\cite{allen2009wp,allen2011china}。

\begin{example}{江南一家的收支}\label{ex:divergence-welfare}
Allen 等人用 Pomeranz 自己的数字算长江三角洲的农家：“Pomeranz 估计，一个全年受雇的男性农业劳动者大约能挣 12 两银子”；按 1745—1754 年的平均价格，“维持一家人‘仅够糊口’的费用是 22.59 两”\cite{allen2009wp}。（1）这个劳动者的福利比率是多少？（2）妻子一年纺纱织布 200 天，“Li 与 Pomeranz 都认为这大约是可能的最多天数”，能挣 14.61 两。夫妻合起来的福利比率是多少？（3）按 Allen 等人的做法，家庭预算是一个成年人的篮子乘以三再加 5\%。一个成年人一年的篮子值多少两？（4）伦敦与阿姆斯特丹的劳动者挣的是糊口费用的四倍。要达到这个水平，这个江南家庭的收入要增加到多少倍？
\solution （1）$12/22.59=0.53$：他只挣到糊口所需的 53\%，“勉强养得活自己，更不用说妻子与孩子了”\cite{allen2009wp}。（2）$12+14.61=26.61$ 两，$26.61/22.59=1.18$；Allen 等人说，“只要不出什么意外，一家人靠这些可以活下去”\cite{allen2009wp}。（3）$22.59/(3\times1.05)=7.17$ 两。（4）$4/1.18=3.4$ 倍。这里拿农家（男性农业劳动者加上妻子的收入）比城市的非熟练劳动者，口径不完全相同；Allen 等人的结论是，即使按 Pomeranz 自己的数字，江南的生活水平“也远远落后于伦敦或阿姆斯特丹”\cite{allen2009wp}。
\end{example}

\subsubsection{产出：人均 GDP}
工资只是收入的一部分。Broadberry、Guan 与 Li 从产出一方重建了中国 980—1850 年的国民核算\cite{bgl2017,bgl2018}。他们的结论是：“中国在北宋时期的生活水平领先世界，但到 1300 年已经落后于意大利”；“到 1750 年，差距已经大到不能用中国内部的地区差别来弥合，大分流在工业革命以前就已经开始了”\cite{bgl2017}。按他们的表 6（1990 年国际元），中国的人均 GDP 1090 年是英国的 116.4\%，1400 年 94.7\%，1500 年与 1600 年都是 77.0\%，1700 年 70.6\%，1750 年 42.5\%，1800 年 29.5\%，1840 年 23.8\%，1850 年 20.0\%\cite{bgl2017}。中国的下降不只是相对的：“到 1840 年降到 980 年水平的约 70\%”\cite{bgl2017}。图~\ref{fig:div_wages}(b) 画出了他们的表 8 里六个经济体 1600—1850 年的人均 GDP。

江南怎么办？按他们的转述，Li 与 van Zanden（2012）发现，1820 年代长江三角洲的人均收入是荷兰的 53.8\%；Broadberry 等人据此推算，下长江地区 1820 年的人均 GDP 约为 1,050 元（1990 年国际元），“比全中国约高 75\%”\cite{bgl2017}。往前推，“长江三角洲的人均 GDP 晚至 1700 年仍是欧洲领先国家水平的 93\%”；“到 1750 年，荷兰的人均 GDP 比长江三角洲高 42\%，到 1770 年领先 90\%”\cite{bgl2017}。这与 Pomeranz 的说法只差在时间：1700 年江南确实与欧洲最富的地区相差不多，可是到 1750 年已经落后了。他们说，这似乎是一种正在形成的“新共识”，既来自 Pomeranz（2011）这样的加州学派作者，也来自用别的数量指标的经济史学者\cite{bgl2017}。Li 与 van Zanden 的论文本书没有取得，只据 Broadberry 等人的转述。

Broadberry 等人的数字是 1990 年国际元，Maddison 数据库 2023 年版是 2011 年国际元，两者的水平不能直接比较；数据库里中国的序列本身就来自他们的研究（经 Xu 等人补充），所以两者的方向相同\cite{mpd2023,bgl2018}。还有一个下限的检验。他们给出 1850 年中国的人均 GDP 是 600 元；另一种算法“得到 472 元，这已经很接近 400 元的仅够糊口的水平，因而是一个有效的下限”\cite{bgl2017}。

\subsubsection{户部的民数可信吗}
\chref{ch:smith}用《清史稿》的户部民数算过，1764—1795 年人口平均每年增长 1.19\%，同时提醒，登记的办法可能变了。Broadberry 等人的人口序列也以官方记录为基础：“在中国，人口数据在这里所涉及的整个时期都有系统的记录，与征税有关。”他们说，“一般认为，官员是能够准确记录人口数目的”，而且“官员有很强的动机正确地估计人口”\cite{bgl2017}。按他们的数字，清代人口的年增长率“提高到每年 0.70\%”\cite{bgl2017}。1764—1795 年登记民数的 1.19\% 比它快：或者 18 世纪后期确实是增长最快的时候，或者登记的办法变了，两种可能都在，本章的数据分不开。回到斯密的判断：清代中国的人口在增长，人均产出与实际工资却在下降。按人均算，斯密说的停滞还说轻了。

\begin{case}{江南与英格兰 · 1700—1850 · 两个核心地区}
Pomeranz 的贡献是改变了比较的单位：不拿中国比欧洲，而拿江南比英格兰。这一比较的三方各有最强的论据。Pomeranz 一方：晚至 1750 年，两地的相似之处多于不同之处，Lal 也承认他的论证“相当令人信服”\cite{lal2000}。Allen 等人一方：即使按 Pomeranz 自己的数字，江南农家一对夫妻的收入只是糊口费用的 1.18 倍，伦敦与阿姆斯特丹的劳动者是四倍（例~\ref{ex:divergence-welfare}）；苏州与广州的福利比率与北京相差无几\cite{allen2009wp}。Broadberry、Guan 与 Li 一方：1700 年江南还是欧洲领先国家的 93\%，1750 年荷兰已经领先 42\%，1770 年领先 90\%\cite{bgl2017}。三方的差别，一半在数据（工资还是产出、哪个城市、哪一年），一半在时间（1700 年还是 1750 年）。今天的共识大致是：分流在工业革命以前就开始了，但开始得比旧说法晚，而且起初只是英格兰与低地国家同其他所有地区的分流，包括欧洲的其他地区。
\end{case}

\fig{div_wages}{生活水平的分流，1600—1913 年。(a) 看伦敦与其他城市的差距什么时候拉开：八个城市非熟练男工的福利比率（全年 250 天的收入除以一家人“仅够糊口”的费用，1 即糊口），每十年平均，没有数据的年份不计；中国的城市用实线，其他用浅色线。(b) 看 1700 年以后的差距：Broadberry、Guan 与 Li 表 8 中六个经济体的人均 GDP，按 1990 年国际元计，纵轴为对数刻度，与图~\ref{fig:div_world} 的 2011 年国际元不能直接比较。数据：(a) Allen 等（2011）的 GPIH 工作簿；(b) Broadberry、Guan 与 Li（2017）表 8；十年平均由本书计算。}{fig:div_wages}

\section{为什么是英国}\label{sec:divergence-why}

这一节回答：工业化为什么先在英国发生？各种解释的最强论据与最大弱点各是什么？表~\ref{tab:divergence-why} 把它们放在一起。

\subsubsection{高工资与便宜的能源（Allen）}
Allen（2009）的回答，按 Leunig 书评引的原文，是英国的“工资高得出奇，能源便宜得出奇”\cite{allen2009,leunig2011}。在 2018 年的一篇答辩里，Allen 自己概括了这个论题：“工业革命在英国被触发，是因为这个国家异常高的工资与便宜的能源价格”；“英国的要素价格使使用——因而发明——提高劳动生产率的机器技术有利可图，这种技术靠的是增加每个工人使用的资本与能源”；“这种不寻常的价格结构，可以归因于英国的帝国政策造成的贸易繁荣”\cite{allen2018wheels}。机器在英国是“独一无二地有利可图”的\cite{leunig2011}，在工资低、煤贵的地方则不划算，所以没有人去发明。这正是马尔萨斯说的中国：“劳动如此充裕，没有人费心去节省它”。

Leunig 认为这本书的大部分论证站得住：第二章“提出了 18 世纪英国工人按国际标准是富裕的证据”，“除了 Greg Clark，所有人都会大体接受这一章”；第三章把农业的标准故事倒了过来，不是农业生产率提高使人离开农村去城市，“相反，城市的高工资使人们离开，这又导致生产率上升”\cite{leunig2011}。第十章是 Allen 对 Mokyr 的看法，“承认启蒙的经济是必要的，但用欧洲的证据论证它是不充分的”\cite{leunig2011}。

棉纺织业是这个论题的主要例子。Broadberry 与 Gupta 说，“国内生产者没有试图模仿印度劳动密集的生产方法，那在英国的工资水平下在经济上不可能行得通”；到 18 世纪中叶，“由英国比印度高得多的工资所推动的、对节省劳动的发明的探索，这时已经开始了”\cite{bg2005}。早在 1600 年，“印度的非熟练劳动者挣的只比英格兰非熟练工资的 20\% 多一点”；1701 年的一位作者写道，“在英格兰要花一先令的劳动，在印度两便士就能得到”\cite{bg2005}。生活水平不像白银工资看起来那么悬殊，因为“印度的谷物价格也相对便宜”；可是“同印度生产者竞争的英国厂商必须按白银工资考虑问题，因为它们必须按世界的白银价格在世界市场上出售”\cite{bg2005}。

\subsubsection{纺纱工的工资之争}
这个论题最尖锐的检验是纺纱。按 Humphries 与 Schneider 对 Muldrew 的转述，手工纺纱“到 1750 年雇用了一百多万妇女与儿童”\cite{humphries2019spin}。如果 Allen 对，纺纱工的工资在珍妮纺纱机等发明之前应当很高。Humphries 与 Schneider 收集了大量纺纱工的工资记录，结论相反：“纺纱这个工业化的主导部门，并不存在这样的‘高工资经济’”，“机械化的动机必须到别处去找”\cite{humphries2019spin}。工资“确实在缓慢上升，但速度只有每年 0.3\%”；“在 1760 年代与 1770 年代的纺纱发明之前，纺纱工的名义工资、实际工资与相对工资都没有大的跳跃”；“即使按当时女性劳动市场的标准，纺纱的报酬也相对较低”\cite{humphries2019spin}。在 18 世纪后期的三项纺纱发明之前，纺纱工的工资“没有出现高峰”，也就是没有“冒烟的枪”\cite{humphries2019spin}。他们的附录 I 给出了每十年的平均：1750—1759 年有 703 个观测，日工资 3.68 便士，是 Allen 的“体面”（respectability）篮子一个成年人每天费用 3.07 便士的 1.20 倍，是女性零工日工资 6.31 便士的 0.58；1760—1769 年只有 26 个观测，3.19 便士，0.97 倍与 0.49；1770—1779 年 60 个观测，4.47 便士，1.20 倍与 0.72\cite{humphries2019spin}。这里的福利比率，按他们的说明，是“用名义工资除以 Allen 的‘体面’篮子里一个成年人的生活费用”\cite{humphries2019spin}，与 Allen 等人比较城市的家庭福利比率口径不同。在他们看来，“制造商作为一个买方垄断的卡特尔行事”；“克服低生产率与纱线质量的不稳定，利用就业不足的妇女与儿童的低工资，这些也许才是修补匠与发明家的动机”\cite{humphries2019spin}。

Allen 的答辩题为《空转他们的纺车》。他认为，Humphries 与 Schneider “不知道纺纱工工作了多少小时，所以这些数据不能确定她们每小时的工资（这才是相关的变量）是高还是低”；而且，“这些信息主要来自学校与慈善项目，项目的参加者正是因为没有生产力才被选中的”\cite{allen2018wheels}。按他的估计，“典型的纺纱工只工作全时的 40\%，只挣到全时收入的 40\%”\cite{allen2018wheels}。他用 Guest（1823）给出的 1760 年纺纱与织布的成本来算纺纱工的报酬：纺纱“每磅花 9 便士，清棉与梳棉合起来又是 9 便士”；“纺一磅棉花需要 1.09 天的劳动，纺纱的全部工序，包括准备与梳棉，共需 2.18 天”\cite{allen2018wheels}。他对比的是 Humphries 与 Weisdorf 的数字：“1750 年代与 1760 年代受雇为仆人的妇女隐含的平均日收入是每天 8.275 便士”\cite{allen2018wheels}。他也换了一个角度：“如果工业革命发生在印度，Humphries 与 Schneider 关于低工资导致机械化的看法就值得认真考虑。但它没有发生在印度。”\cite{allen2018wheels}他的结论是：“高工资假说并没有因为 Humphries 与 Schneider 的批评而受损。”\cite{allen2018wheels}

Humphries 与 Schneider 的回应“为我们的数据集增加了 1500 多个新的观测值”，结论不变：纺纱工的工资“即使同其他女工相比也很低，其变化轨迹也不能解释珍妮纺纱机的发明与推广”\cite{humphries2019wheel}。他们“质疑这样一个奇怪的假定：纺纱工虽然贫困，却选择放弃挣钱的机会而宁要闲暇”；还回敬了 Allen 的题目：“在原地空转纺车的是 Allen 教授，不是我们。”\cite{humphries2019wheel}

\begin{example}{纺纱工一天挣多少}\label{ex:divergence-spinners}
（1）按 Allen 引的 Guest 的数字，纺一磅棉花付 9 便士，需要 1.09 天。纺纱工一个全天挣多少？包括清棉与梳棉的全部工序共付 18 便士、需要 2.18 天，一天又是多少？与 Humphries 与 Weisdorf 的女仆日收入 8.275 便士相比如何？（2）Humphries 与 Schneider 的 1750 年代平均日工资是 3.68 便士。它是（1）的多少？如果像 Allen 说的，典型的纺纱工只工作全时的 40\%，按（1）算，她一个“工作日”的收入是多少？（3）两方的分歧在哪里？
\solution （1）$9/1.09=8.26$ 便士；$18/2.18=8.26$ 便士，两种算法相同，与 8.275 便士几乎一样。按每个全天的报酬算，纺纱工挣得与女仆一样多。（2）$3.68/8.26=0.45$。$0.4\times8.26=3.30$ 便士，与 3.68 便士相差不多。（3）分歧不在算术，而在“一天”是什么：Humphries 与 Schneider 记录的是纺纱工实际拿到的日收入，Allen 认为那只是部分时间的收入，相关的是每小时的报酬。如果纺纱工多数只做部分时间，日收入低与小时报酬高可以同时成立；Humphries 与 Schneider 反对的正是这个假定：贫困的纺纱工为什么会放弃挣钱的机会？数据回答不了谁对，因为没有人记下纺纱工每天工作了多少小时。
\end{example}

\subsubsection{煤与殖民地（Pomeranz）}
按 Lal 的概括，Pomeranz 的回答分两部分\cite{pomeranz2000,lal2000}。第一，“为英国工业革命提供燃料的煤，在地理上对 18 世纪中国南方的核心地区并不那么容易得到，因为它集中在西北”；回头看，“把长江三角洲与西北的煤矿连起来，回报是巨大的，但这些回报事前是看不见的”\cite{lal2000}。第二，“Pomeranz 对大分流原因的回答的第二部分，是欧洲对新大陆的发现与开发——部分是通过贸易”\cite{lal2000}。Lal 也说，这“无疑扩展了欧洲的土地边疆”\cite{lal2000}；Broadberry、Guan 与 Li 把 Pomeranz 的这两个因素概括为煤与“幽灵土地”（ghost acres），即欧洲在本土以外用到的土地\cite{bgl2017}。中国也有一个人口稀少的边疆：Pomeranz 说，东南亚“像接触以后的新大陆一样，人口稀少，能够供给大量‘国内’需要的土地密集型资源；中国人大量去了那里，但东南亚从来没有成为沿海中国的新大陆”，因为中国人是“在没有国家支持的情况下在东南亚立足”的\cite{lal2000}。用式~\eqref{eq:divergence-escape} 的话说，煤与新大陆降低了土地在英国生产里的份额 $\alpha$。

Lal 不接受这个回答。关于煤，他说，11 世纪中国西北“由 Hartwell 记录的煤铁联合体的惊人发展”，被 12 世纪的入侵者毁掉了；中国在北宋时“几乎具备了这场革命的所有要素，可它仍然没有发生”\cite{lal2000}。（Hartwell 的数字后来受到质疑：他说 1078 年中国产铁多达 15 万吨，按 Broadberry 等人采用的 Liu（1993）的估计，当年的最高产量约为 1.35 万吨\cite{bgl2017}。）印度古代的核心地区“在恒河平原的东部，即今天的比哈尔”，“我们现在知道，这个地区也有印度的煤炭储藏”，“尽管如此，没有人认为印度人本可以发展出煤与蒸汽的工业革命”\cite{lal2000}。关于新大陆，Pomeranz 自己“承认从新大陆流入欧洲的土地密集型产品，增量不可能很大”\cite{lal2000}。Lal 的结论是：“所以，正如韦伯以及近来的 Landes 都坚持的，我们又回到了文化。”\cite{lal2000}这里的韦伯是马克斯·韦伯（Max Weber，1864—1920）。

\subsubsection{知识（Mokyr）}
Mokyr（2009）的解释是\textbf{工业启蒙}：英国的技术进步来自“我们与启蒙运动联系在一起的那套不断变化的信念”\cite{mokyr2009,harley2010}。按 Harley 的概括，这是“一个社会，其中占主导的是这样一种心态：相信通过观察与实验获得的有用知识可以带来进步”；“持续的变化取决于持续的技术进步，而技术进步又取决于知识及其发展的过程”\cite{harley2010}。用式~\eqref{eq:divergence-escape} 的话说，这是 $g$ 的来源，而且不是一次性的提高，而是让 $g$ 持续下去的机制。Harley 的批评有两点。第一，“所谓工业启蒙，似乎至少是在描述成功的变化是什么样子，而不是它的原因”；第二，这本书“没有完全说服人：为什么一场英国特有的工业启蒙，应当在理解 18 世纪现代经济增长的出现时占首要的位置”\cite{harley2010}。Allen 的看法前面说过：启蒙是必要的，但不充分，欧洲许多地方都有启蒙，只有英国有高工资与便宜的煤\cite{leunig2011}。

\subsubsection{文化与人口（Clark）}
Clark（2007）认为，此前也有过技术改进，但“问题在于，这样的改进太零星、太有限了”；“要得到持续的经济增长，我们需要持续的全要素生产率增长”\cite{clark2007,margo2008}。他的解释是，漫长的前工业时代“在文化上，或许也在遗传上”塑造了人（Margo 引的原文），Clark 指的是西方关键行为的发展，“尤其是耐心”\cite{clark2007,margo2008}。按 Margo 的概括，“第六章根据遗嘱提出证据，表明生育上的成功（男人有多少个孩子）与去世时的财富正相关”：财富如果来自耐心与勤劳这些“中产阶级的”价值，这些行为就可以代代相传，Clark 相信确实如此，“只等着技术的持续进步一旦到来就加以利用”\cite{margo2008}。为什么不是中国？“Clark 的回答是，中国等地的人口增长更快（更马尔萨斯），富人生育上的成功更少”；他也不接受制度的解释，“认为所有相关的制度在工业革命之前很久就已经在英格兰具备了”\cite{margo2008}。Margo 的批评很直接：“即使只看表面，我该相信这些数据吗？”“让良好的行为代代相传的，究竟是什么机制？”他的总评是：“我认为，任何单一的故事都不太可能解释国家的财富。”\cite{margo2008}

\subsubsection{制度（Acemoglu、Johnson 与 Robinson）}
\chref{ch:islamic}讨论马格里布商人时留下了一个问题：靠关系网络执行的契约，长期看离不开公共的法律制度。Acemoglu、Johnson 与 Robinson（2002）把制度与贸易连了起来\cite{ajr2002}。他们发现，“1500—1850 年（西）欧的兴起，主要可以由有通向大西洋通道的欧洲国家的增长来解释，尤其是那些从事殖民与远洋长途贸易的国家”；机制是，“1500 年以后同新大陆、非洲与亚洲贸易的增长，壮大了商业资产阶级中的新的部分，使这些群体能够要求、得到并维持保护其产权的制度变革”\cite{ajr2002}。条件是原有的制度：“最重要的制度变化，以及因此最大的经济收益，出现在原有制度对君主有某些制约、尤其是限制了君主对海外贸易控制的国家。”\cite{ajr2002}所以，“欧洲的兴起，主要是晚期中世纪的制度与‘大西洋贸易’提供的经济机会相互作用所推动的资本主义发展的结果”\cite{ajr2002}。这与\chref{ch:finrevolution}讲的 1688 年以后英国国债的可信，是同一个故事的两面。Broadberry、Guan 与 Li 从中国的数据得出类似的结论：分流在 1750 年以前就开始了，“这表明大分流有深厚的制度根源，而不是因为煤或‘幽灵土地’之类的因素突然出现的”\cite{bgl2017}。

\begin{table}[htbp]\centering\small
\caption{为什么是英国：五种解释}\label{tab:divergence-why}
\begin{tabularx}{\linewidth}{@{}>{\raggedright\arraybackslash}p{0.13\linewidth}>{\raggedright\arraybackslash}p{0.14\linewidth}YYY@{}}
\toprule
解释 & 在式~\eqref{eq:divergence-escape} 里 & 机制 & 最强的论据 & 最大的弱点 \\\midrule
高工资与便宜的能源（Allen） & $g$ 提高 & 工资高、煤便宜，节省劳动的机器在英国有利可图，于是被发明 & 1770 年英国棉纺织业的工资是印度的 6 倍；伦敦的福利比率是北京的 3 倍 & 手工纺纱工的工资低，发明前没有上升（Humphries 与 Schneider） \\
煤与殖民地（Pomeranz） & $\alpha$ 下降 & 煤代替木柴，新大陆的土地补上本国的土地 & 江南没有近便的煤，东南亚没有成为中国的新大陆 & 宋代与印度的煤；新大陆的增量不大（Lal） \\
知识（Mokyr） & $g$ 持续 & 工业启蒙使有用的知识不断积累 & 技术进步在 1800 年以后没有停下来 & 描述多于原因；不一定是英国特有的（Harley） \\
文化与人口（Clark） & $g$ 与 $\beta$ & 富人的后代多，中产阶级的习惯扩散 & 遗嘱里财富与子女数正相关 & 数据与传递机制可疑（Margo） \\
制度（Acemoglu 等） & $g$ 与 $\alpha$ & 大西洋贸易壮大了商人，商人约束了君主 & 有大西洋通道、君主受制约的国家增长最快 & 制度何时成形、为什么在英格兰（Clark：早已具备） \\
\bottomrule
\end{tabularx}
\datasource{本节所引各书与论文；Pomeranz、Allen、Mokyr 与 Clark 四本书的论点据 EH.Net 的书评。第二栏是本书的归类。}
\end{table}

\begin{compare}{高工资解释与资源解释}{高工资（Allen）}{资源（Pomeranz）}
\cmprow{关键的价格}{劳动相对于资本与能源贵}{土地密集的产品与燃料稀缺}
\cmprow{煤的作用}{便宜的煤使蒸汽机与机器划算}{煤就在需要的地方，是运气}
\cmprow{殖民地的作用}{帝国的贸易抬高了英国的工资}{新大陆提供了“幽灵土地”}
\cmprow{为什么不是江南}{工资低，节省劳动不划算}{煤在西北，东南亚不是新大陆}
\cmprow{为什么不是印度}{工资更低（只有英国的 20\% 多一点）}{没有直接回答；Lal：比哈尔也有煤}
\cmprow{分流的开始}{17—18 世纪英格兰与低地国家的高工资}{1800 年前后，煤与新大陆起作用以后}
\end{compare}

\begin{case}{印度 · 1600—1850 · 棉纺织业的兴衰}
Broadberry 与 Gupta 说，17 世纪起东印度公司把印度棉布运进英国，“为英国制造商开辟了新的机会，即进口替代与转口替代的策略”；到 19 世纪初，英国却“已经成为世界上最重要的棉纺织品生产国，支配着世界的出口市场，甚至向印度出口”\cite{bg2005}。起初，英国用法律挡住印度的棉布：“从 1701 年起，从印度进口的印花布与其他某些种类的棉布被禁止”；“1701 年的《印花布法》仍然允许从印度进口白棉布在英国印花，直到 1721 年的进一步立法禁止了这种进口，除非是为了转口”；这项法律经过修正，“一直有效到 1774 年”\cite{bg2005}。Broadberry 与 Gupta 认为，“保护对英国棉纺织业的作用似乎是中性的，《印花布法》被规避了”；而且，“如果《印花布法》成功了，英国棉纺织业也许会像美国的同行一样，永远不能在世界市场上竞争”\cite{bg2005}。

决定竞争力的是单位劳动成本，即工资除以劳动生产率。以印度 $=100$，记英国的白银工资为 $W/W^{*}$、劳动生产率为 $y/y^{*}$，原棉价格为 $C/C^{*}$，离岸价格为 $P/P^{*}$，全要素生产率为 $A/A^{*}$：
\begin{keyeq}
\frac{ULC}{ULC^{*}}=\frac{W/W^{*}}{y/y^{*}},\qquad \frac{A}{A^{*}}=\frac{(W/W^{*})^{\theta}(C/C^{*})^{1-\theta}}{P/P^{*}} .\label{eq:divergence-ulc}
\end{keyeq}
第二式是成本对偶的全要素生产率，$\theta$ 是工资在成本里的权重；Broadberry 与 Gupta 的表 13 取 $\theta=0.5$\cite{bg2005}。按他们的表 7，兰开夏棉纺织工人的周薪，约 1680 年是 6 先令，印度是 1 先令 6 便士，英国是印度的 400\%；约 1770 年 9 先令对 1 先令 6 便士，600\%；约 1790 年 13 先令 3 便士对 2 先令，663\%；约 1820 年 10 先令 4 便士对 2 先令，517\%\cite{bg2005}。“1770 年，英国的工资约为印度的 6 倍，但英国的劳动生产率只比印度高两到三倍”：按表 10，英国的劳动生产率是印度的 250\%，单位劳动成本是 240\%，英国贵得多\cite{bg2005}。此后机器出现了。Cuenca Esteban 估计，棉纺织品出口的实际价格 1770/4—1820/4 年间下降到 1/6.53；Broadberry 与 Gupta 据此说，“到 1820 年，由于英国的生产率提高到 6.53 倍，并假定印度停滞，英国的生产率上升到印度水平的 1625\%”，“由于这时英国的工资已经回落到印度水平的 517\% 左右，英国的单位劳动成本已经低于印度”：$517/1625=32\%$\cite{bg2005}。

结果在印度很快就看得见。孟加拉输往英国的棉布，1790—1799 年每年 78.7 万匹，1830—1839 年只有 6,000 匹；全印度输往英国的，从 220 万匹降到 27.1 万匹\cite{bg2005}。“1818 年，达卡的商馆关闭了。促使董事会在 1818 年放弃它的，不是运费，也不是保护性关税，而是英国同类棉制品价格的大幅下降”，孟加拉的细布越来越失去市场（Broadberry 与 Gupta 引 Mitra）\cite{bg2005}。印度本身的市场守得更久，因为运输成本“一直到 1860 年代都给了印度生产者在本国市场上的优势”\cite{bg2005}。按表 4，印度的棉布消费 1831—1835 年是 3.75 亿磅，英国货占 3.9\%；1856—1860 年 4.55 亿磅，占 35.3\%；1880—1881 年 6 亿磅，占 58.4\%\cite{bg2005}。按 Twomey 的估计，印度虽然从 1830 年前后成为棉布的净进口国，供应本国市场的手工业生产“只是在 1850 年以后才下降”\cite{bg2005}。英国机器打败印度手工业的速度，取决于两国的工资、机器的生产率、原棉的价格与运费；印度在第一次全球化中的位置见\chref{ch:periphery}。
\end{case}

\begin{example}{单位劳动成本与全要素生产率}\label{ex:divergence-ulc}
（1）按式~\eqref{eq:divergence-ulc} 与表 10 的数字，算 1770 年与 1820 年英国的单位劳动成本（印度 $=100$）。（2）1770 年英国的劳动生产率是印度的 250\%，此后提高到 6.53 倍，1820 年应是多少？Broadberry 与 Gupta 印的是 1625\%，差在哪里？（3）表 13 的原棉价格之比，1770 年是 320，1820 年是 127；离岸价格之比是 200 与 43。取 $\theta=0.5$，按式~\eqref{eq:divergence-ulc} 的第二式，两年的全要素生产率之比各是多少？（4）约 1680 年，工资之比 400，原棉 182，离岸价格 200。1680—1820 年，英国全要素生产率对印度的优势每年扩大多少？
\solution （1）$600/250=240\%$；$517/1625=31.8\%$，表 10 印作 32。（2）$250\times6.53=1{,}632.5\%$；1625 是 $250\times6.5$，大概是取了整数倍。（3）1770 年：$\sqrt{600\times320}=438$，$438/200=2.19$，即 219；1820 年：$\sqrt{517\times127}=256.2$，$256.2/43=5.96$，即 596。与表 13 一致。（4）1680 年：$\sqrt{400\times182}=270$，$270/200=1.35$，即 135。$\ln(596/135)/140=1.06\%$。这正是 Broadberry 与 Gupta 脚注 8 说的每年 1.06\%，脚注把时期印作“1680—1770”，按表算应是 1680—1820 年；1680—1770 年只有 $\ln(219/135)/90=0.54\%$。
\end{example}

\fig{div_cotton}{英国棉纺织业的崛起。(a) 看原棉消费量涨了多少倍：英国每年的原棉消费量（每十年平均，百万磅），纵轴为对数刻度。(b) 看英国的成本优势从哪里来：英国对印度的工资、原棉价格、两者的加权成本、离岸价格与全要素生产率（印度 $=100$），约 1680、1770、1790、1820 年，纵轴为对数刻度；虚线是 100，即与印度相同。数据：Broadberry 与 Gupta（2005）表 1 与表 13。}{fig:div_cotton}

\section{数据检验与反例}\label{sec:divergence-test}

这一节回答：英格兰 1250—1860 年的数据，支持马尔萨斯模型吗？模型从什么时候开始不再适用？

\subsubsection{收入与人口的边界}
按式~\eqref{eq:medieval-factor}，技术与土地不变时，人均收入对人口的弹性是 $-\alpha$：人口增加 1\%，人均收入下降 $\alpha$\%。Clark 的十年数据给出了英格兰的人口与人均实际国民净收入，也给出了地租与工资收入，可以算出土地的份额 $\alpha=$ 地租 $/($地租 $+$ 工资收入$)$\cite{clark2010}。1250—1630 年，$\alpha$ 平均是 0.235，在 0.163 与 0.314 之间。

图~\ref{fig:div_england}(a) 把每个十年画成一个点。1250—1630 年的 39 个十年排成一条向下的线：用对数人均实际国民净收入对对数人口回归，斜率是 $-0.72$，95\% 置信区间（按 Newey–West 方法，滞后三期）为 $-0.83$ 到 $-0.60$，$R^{2}=0.88$；农业实际工资的斜率更陡，是 $-1.17$\cite{clark2010}。方向与模型一致，幅度却是土地份额的三倍。原因可能有三个。第一，模型假定技术不变，而 1250—1630 年的技术、气候与疾病环境都在变；如果人口多的时期正好也是收成坏的时期，弹性就会被夸大。第二，Clark 的收入从工资与地租算起，再用生活费用指数平减；如果这个指数以食物为主，人口多时粮价高，平减以后的收入就下降得更多。第三，土地与劳动之间的替代弹性可能小于 1，人多时土地的份额本身就会上升；Clark 的数据确实如此，1250—1630 年，土地份额与对数人口的相关系数是 0.53。用 Broadberry 等人的产出数据，1270—1630 年的斜率只有 $-0.28$（95\% 置信区间 $-0.50$ 到 $-0.05$，$R^{2}=0.36$），正好在土地份额附近\cite{broadberry2015,boe2017}。哪一套数据更接近事实，决定了马尔萨斯模型在中世纪的英格兰是“对了方向”还是“完全对”。

1640 年以后，这条线断了。1640—1790 年，斜率变成 $+0.45$；1800—1860 年是 $+0.46$，$R^{2}=0.995$：人口与人均收入一起增长\cite{clark2010}。马尔萨斯的边界在 17 世纪中叶以后开始向外移动，在 19 世纪移动得越来越快。

\subsubsection{全要素生产率}
土地不变时，由 $y=\tilde A(X/L)^{\alpha}$，全要素生产率的增长率是
\begin{equation*}
\Delta\ln\tilde A=\Delta\ln y+\alpha\,\Delta\ln L ,
\end{equation*}
人均收入的增长加上人口增长被土地报酬递减吃掉的部分。逐个十年按前后两个十年的平均 $\alpha$ 计算，再按年平均：1250—1640 年，全要素生产率每年变化 $-0.04\%$，人均收入每年 $-0.06\%$，人口每年 $+0.09\%$；1640—1800 年，全要素生产率每年增长 0.25\%，人均收入 0.18\%，人口 0.32\%；1800—1860 年，全要素生产率 0.80\%，人均收入 0.61\%，人口 1.29\%\cite{clark2010}。土地的份额同时在下降：1250 年代 0.284，1600 年代 0.314，1700 年代 0.252，1800 年代 0.211，1850 年代 0.109，1860 年代 0.095。

把这些数字放进式~\eqref{eq:divergence-escape}。1250—1640 年，$g$ 几乎是零，经济停在马尔萨斯的稳态附近，人口与收入此消彼长。1640—1800 年，$g$ 是正的，可是低于门槛：$\alpha$ 约 0.25，$\alpha n_{\max}=0.69\%$，人口如果全速增长，仍可能吃掉全部进步；实际上人口每年只增长 0.32\%，收入增长了。1800—1860 年，全要素生产率每年 0.80\%，$\alpha$ 平均 0.149，门槛只有 0.41\%：三条出路中的两条（$g$ 提高与 $\alpha$ 下降）同时出现，人口增长再快也吃不掉进步。第三条（$\beta$ 下降）要到 19 世纪后期的生育率下降才出现。

\begin{example}{全要素生产率的核算}\label{ex:divergence-tfp}
（1）1640—1800 年，英格兰的人均实际收入每年增长 0.177\%，人口每年增长 0.323\%，土地份额平均 0.25。全要素生产率每年增长多少？其中多大比例表现为人均收入的增长？（2）1800—1860 年，人均收入每年增长 0.609\%，人口 1.290\%，土地份额平均 0.15。再算一遍。（3）1700—1800 年，Clark 的人口增加到 1.651 倍。如果技术不变、$\alpha=0.25$，人均收入应变为多少？实际上，人均实际国民净收入从 58.97 升到 69.63，农业实际工资从 70.09 降到 67.34。这说明什么？
\solution （1）$0.177+0.25\times0.323=0.258\%$；$0.177/0.258=69\%$。（逐个十年按变化的 $\alpha$ 算是 0.253\%。）（2）$0.609+0.15\times1.290=0.803\%$；$0.609/0.803=76\%$。（3）$1.651^{-0.25}=0.882$，人均收入应下降 11.8\%。实际上人均收入增加了 18.1\%，全要素生产率的进步不只抵消了人口的增长。可是农业工资下降了 3.9\%，建筑业的非熟练工与工匠的实际工资也分别从 58.17 降到 56.87、从 60.83 降到 55.90。马尔萨斯说国家财富的增长没有改善劳动贫民的处境，到 1800 年为止，按工资看是对的；增加的收入大部分没有落到工资上。1800—1850 年，农业、非熟练工与工匠的实际工资才分别增加 45.9\%、62.2\% 与 63.2\%。
\end{example}

\fig{div_england}{英格兰的马尔萨斯边界，1250—1860 年。(a) 看三团点的斜率：每个十年一个点，横轴是人口（百万），纵轴是 Clark 的人均实际国民净收入指数（1860 年代约为 100），两轴都是对数刻度；直线是 1250—1630 年的回归线，斜率 $-0.72$。(b) 看两套估计在哪里分开：Clark 的人均实际国民净收入（实线）与 Broadberry 等人的人均实际 GDP（褐色虚线，十年平均，1860 年代 $=100$）。数据：Clark 的十年数据（Clark 2010 及其工作簿）；英格兰银行《千年宏观数据》表 A21（Broadberry 等 2015）；回归与十年平均由本书计算。}{fig:div_england}

\begin{pitfall}[马尔萨斯在 1798 年就错了，所以他的模型没有用]
马尔萨斯写书的时候，英格兰正在离开他描述的经济，这是事实。可是他的模型解释了此前几百年的英格兰：1250—1630 年，人均收入与人口此消彼长，弹性是负的；全要素生产率几乎不增长。连他对穷人的判断，到 1800 年为止也是对的：1700—1800 年人均收入增加了，农业工资却没有。式~\eqref{eq:divergence-escape} 说明，模型并不排斥增长：只要全要素生产率增长快过土地份额乘以人口增长的上限，收入就会增长；模型错的是一个隐含的假定，即技术进步永远像中世纪那样慢、土地永远那样重要。
\end{pitfall}

\begin{pitfall}[分流是 1800 年以后才开始的]
按工资，伦敦与阿姆斯特丹的劳动者在 18 世纪中叶已经比北京、苏州与广州的劳动者富得多，福利比率是后者的三倍左右（2.9—3.6 倍）；按产出，荷兰的人均 GDP 1750 年已经比长江三角洲高 42\%。工业革命扩大了差距，不是创造了差距。另一方面，分流起初并不是“欧洲与亚洲”的分流：1750 年代米兰的福利比率（1.34）与北京（1.35）几乎相同，莱比锡（1.65）也不到伦敦的一半。最早分开的是英格兰与低地国家同其他所有地区。
\end{pitfall}

\begin{deeper}[版本与史料的出入]
\textbf{四本书。}Pomeranz（2000）、Allen（2009）、Mokyr（2009）与 Clark（2007）都没能取得，本章据 EH.Net 的四篇书评：Lal（2000）、Leunig（2011）、Harley（2010）与 Margo（2008）。凡引号里的话，或是书评人的概括，或是书评引的原文，后者注明页码（Allen 第 2 页、Mokyr 第 478 页、Clark 第 8 页、Pomeranz 第 200 页）。Harley 书评的标题栏把 Mokyr 的书印作“Princeton: Princeton University Press, 2010”，同时给出的 ISBN 978-0-300-12455-2 却是耶鲁大学出版社的号码；本章按通常的引法记为耶鲁大学出版社 2009 年版，没有核对原书\cite{harley2010}。

\textbf{工作论文与正式发表的版本。}Allen 等人（2009）是 2011 年《经济史评论》论文的工作论文版本；Broadberry、Guan 与 Li（2017）是 2018 年《经济史杂志》论文的工作论文版本；Broadberry 与 Gupta（2005）是 2009 年《经济史评论》论文的会议版本；Humphries 与 Schneider 的论文据作者在牛津大学档案库里的录用稿（无印刷页码），他们 2019 年的回应与 Allen 2018 年的答辩据牛津的工作论文；Acemoglu、Johnson 与 Robinson（2002）是 2005 年《美国经济评论》论文的工作论文版本（© 2002 by Daron Acemoglu, Simon Johnson and James Robinson，短引）。本章的页码与表号都按这些版本。

\textbf{Broadberry 与 Gupta 的表。}表 10 的生产率 1625 是 $250\times6.5$，不是正文说的 6.53 倍（例~\ref{ex:divergence-ulc}）；脚注 8 的时期“1680—1770”按表 13 应是 1680—1820 年。式 (1) 的说明写的是“星号表示作为计价标准的国家，这里是英国”，式 (2) 却说“星号表示外国，这里取为印度”；表 10 与表 13 都以印度 $=100$，本章按表的方向，把星号统一为印度\cite{bg2005}。表 13 的离岸价格之比，1770 年以前据 Chaudhuri 的东印度公司加价率，“我们假定最初的离岸价格之比是 200”\cite{bg2005}，此后据 Cuenca Esteban 的出口价格指数推算，假定印度的价格不变；全要素生产率的数字因此依赖这些假定。

\textbf{Humphries 与 Schneider 的附录。}1760—1769 年的福利比率印作 0.97，按印出的日工资与篮子费用算是 $3.19/3.31=0.96$，大概是用未四舍五入的数算的；这个十年只有 26 个观测\cite{humphries2019spin}。

\textbf{两种国际元。}Broadberry、Guan 与 Li 的数字按 1990 年国际元（Maddison 原来的单位），Maddison 数据库 2023 年版按 2011 年国际元。本章只在同一个来源内部比较水平：中国对英国的比率用前者，六个经济体的长序列用后者。

\textbf{没有取得的。}Li 与 van Zanden（2012）关于长江三角洲与荷兰的比较，本章据 Broadberry、Guan 与 Li 的转述；Broadberry 等人关于英格兰的书只用了英格兰银行转载的数据；Perkins（1969）据 Broadberry 等人的转述。
\end{deeper}

\section{今天还剩什么}\label{sec:divergence-legacy}

大分流之争没有结论，它提出的问题贯穿本书后面的许多章。

\textbf{增长从哪里来。}式~\eqref{eq:divergence-escape} 里，$\alpha\to0$、$\beta\to0$ 时，收入每代增长 $g$，土地从模型里消失，剩下的问题是 $g$ 从哪里来。索洛模型把增长分成资本积累与技术进步，技术进步是外生的（\chref{ch:solow}）；内生增长理论问的正是 Mokyr 的问题：知识怎样积累（\chref{ch:endogenous}）。

\textbf{人口转变。}第三条出路，$\beta$ 下降，要到 19 世纪后期才出现。今天许多国家面对的是相反的问题：生育率低于更替水平，人口减少，劳动变得稀缺（\chref{ch:demography}）。

\textbf{发展与赶超。}马尔萨斯说，中国靠劳动者之多、劳动之便宜，可以制造数量极大的制成品向外国出售。两百年后，20 世纪的东亚经济在很大程度上走的是这条路（\chref{ch:eastasia}）。Allen 的问题反过来问：工资低的国家为什么要用发达国家为高工资设计的技术？二元经济与大推进的理论讨论的就是这个问题（\chref{ch:development}）。

\textbf{去工业化。}印度的手工棉纺织业被英国的机器打败，是第一次全球化中边缘国家的经历的一部分（\chref{ch:periphery}）。

\textbf{分配。}1700—1800 年英格兰的人均收入增加了，工资却没有。工业化的成果怎样分配，是下一章马克思的问题（\chref{ch:marx}），也是今天关于不平等的讨论的起点（\chref{ch:heterogeneity}）。

\textbf{数字。}本章的许多争论归结为数字：Clark 与 Broadberry 等人的两套英格兰，工资与产出的两种中国，纺纱工每天工作几小时。宏观数据的可靠性见\chref{ch:datareliability}。

\begin{takeaways}
\begin{itemize}
\item 马尔萨斯（1798）：人口按几何级数增长，生活资料按算术级数增长；预防性抑制与积极性抑制把人口压回来，财富的增长不能改善劳动贫民的处境。式~\eqref{eq:divergence-escape} 把\chref{ch:medieval}的模型推广为 $z_{t+1}=(1-\alpha\beta)z_t+g$：技术进步只把收入推上一个平台 $g/(\alpha\beta)$，此后都变成人口。
\item 三条出路：$g$ 提高（知识、诱致性发明）、$\alpha$ 下降（煤、新大陆、工业）、$\beta$ 下降（晚婚少生）。全要素生产率增长快过 $\alpha n_{\max}$，人口增长再快也吃不掉进步。
\item 起飞很慢：英国的人均 GDP 1700—1820 年每年增长 0.26\%，1820—1870 年 1.13\%；原棉消费从 1700 年代的 110 万磅增加到 1840 年代的 5.263 亿磅。1820 年以前的数字都是估计，Clark 与 Broadberry 等人的英格兰在 1700 年以前相差很大。
\item 分流从什么时候开始：Pomeranz 说 1750 年江南与英格兰相似；Allen 等人的福利比率显示，18 世纪中叶伦敦是北京的三倍，中国的城市与意大利并列最后；Broadberry、Guan 与 Li 的产出显示，1700 年江南还是欧洲领先国家的 93\%，1750 年荷兰领先 42\%，1770 年领先 90\%。
\item 为什么是英国：Allen 的高工资与便宜的能源，Pomeranz 的煤与殖民地，Mokyr 的工业启蒙，Clark 的文化与人口，Acemoglu 等人的大西洋贸易与制度。每一种都有最强的论据，也都有未回答的批评（表~\ref{tab:divergence-why}）；纺纱工的工资之争说明，关键的数字至今没有定论。
\item 印度：1770 年英国棉纺织业的工资是印度的 6 倍，单位劳动成本是 240\%；1820 年生产率是印度的 1625\%，单位劳动成本降到 32\%。孟加拉输往英国的棉布从每年 78.7 万匹降到 6,000 匹。
\item 英格兰的数据：1250—1630 年人均收入对人口的弹性是 $-0.72$，1800—1860 年是 $+0.46$；全要素生产率每年 $-0.04\%$、0.25\%、0.80\%；土地份额从 0.28 降到 0.095。1700—1800 年工资没有上升，1800—1850 年才上升。
\end{itemize}
\end{takeaways}

\begin{selfcheck}
\begin{enumerate}
\item 一个国家有 1,000 万人，人口每 25 年翻一番；食物现在够 1,000 万人吃，每 25 年增加够 1,000 万人吃的一份。100 年后，人口与食物各够多少人？差多少？
\item 设 $\alpha=0.2$，$\beta=0.5$，全要素生产率每年增长 0.4\%，一代 30 年。按式~\eqref{eq:divergence-escape}，平台 $z^{*}$ 是多少？收入是稳态的多少倍？到了平台，人口每年增长多少？离平台的差距多少代减半？
\item 一个劳动者一年挣 30 两，一个成年人一年仅够糊口的篮子值 8 两。按 Allen 等人的做法（三个篮子，再加 5\% 房租），他的福利比率是多少？
\item 英国的工资是外国的 400\%，劳动生产率是 200\%。单位劳动成本是多少？如果工资不变，生产率要提高到外国的多少，英国的单位劳动成本才低于外国？提高到 500\% 时是多少？
\item 某个时期人均收入每年增长 0.3\%，人口每年增长 0.6\%，土地份额 0.2。全要素生产率每年增长多少？其中多大比例表现为人均收入的增长？
\item 技术与土地不变，人口增加到 1.5 倍，$\alpha=0.3$。人均收入变为多少？如果实际上人均收入没有变，全要素生产率提高了多少？
\end{enumerate}
\end{selfcheck}

\begin{answers}
\item 人口 $1{,}000\times2^{4}=16{,}000$ 万；食物够 $1{,}000+4\times1{,}000=5{,}000$ 万人；差 11,000 万人。
\item $g=30\times0.4\%=0.12$，$z^{*}=0.12/(0.2\times0.5)=1.2$，$e^{1.2}=3.32$ 倍；人口每代增长 $\beta z^{*}=0.6$，每年 2\%；差距每代乘以 $1-0.1=0.9$，减半要 $\ln0.5/\ln0.9=6.58$ 代。
\item 家庭预算 $8\times3\times1.05=25.2$ 两，福利比率 $30/25.2=1.19$。
\item $400/200=200\%$；生产率超过 400\%；500\% 时是 $400/500=80\%$。
\item $0.3+0.2\times0.6=0.42\%$；$0.3/0.42=71.4\%$。
\item $1.5^{-0.3}=0.885$，下降 11.5\%；全要素生产率提高了 $1.5^{0.3}-1=12.9\%$。
\end{answers}

英格兰的人均收入在 1800 年以后持续增长，工资却直到 19 世纪中叶才明显上升。马尔萨斯说劳动贫民挤在制造业里、住在狭窄而不卫生的屋子里；半个世纪以后，有人从同一个工厂制度里看到的不是马尔萨斯的人口规律，而是资本积累的规律。下一章讲马克思。
